% Generated draft: requires mathematical and source review.
\documentclass[11pt]{article}
\usepackage[margin=25mm]{geometry}
\usepackage{fontspec}
\setmainfont{Latin Modern Roman}
\usepackage{amsmath,amssymb,mathtools,unicode-math}
\setmathfont{Latin Modern Math}
\usepackage{enumitem,needspace,xurl,hyperref}
\setlength{\parindent}{0pt}
\setlength{\parskip}{5pt}
\setlength{\emergencystretch}{3em}
\newcommand{\N}{\mathbb{N}}
\newcommand{\Z}{\mathbb{Z}}
\newcommand{\Q}{\mathbb{Q}}
\newcommand{\R}{\mathbb{R}}
\newcommand{\C}{\mathbb{C}}
\newcommand{\F}{\mathbb{F}}
\newcommand{\A}{\mathbb{A}}
\newcommand{\PP}{\mathbb{P}}
\newcommand{\E}{\mathbb{E}}
\newcommand{\eps}{\varepsilon}
\newcommand{\dd}{\,\mathrm{d}}
\begin{document}
\section*{UnsolvedMath ID: 11100032}
\textbf{Source problem code:} catalogue entry
\subsection{Definitions and mathematical statement}
\emph{Applications problem 8 (Mark Hovey's list).} Extend the Benson--Carlson--Rickard classification results to connected cocommutative Hopf algebras over a field, with the source's \(A(n)\) as an intended example. The field characteristic, finite-dimensionality, grading/connectedness convention, module category, and exact classification theorem to extend are not supplied.

The stated target is calculation of thick subcategories. A thick subcategory of a triangulated category is a full triangulated subcategory closed under retracts. A thick tensor ideal additionally remains closed under tensoring with arbitrary objects; these notions are different unless an extra argument identifies them. The record does not specify whether it means the stable category of finite modules, a derived category, graded modules, or tensor ideals. A precise extension must fix the category, support space and support theory, and the resulting classification correspondence.

The source reports Hovey--Palmieri reducing the calculation to quasi-elementary Hopf algebras. It does not define quasi-elementary or state the reduction's hypotheses. For group algebras it says quasi-elementary sub-Hopf algebras look like \(k[E]\) for elementary abelian \(p\)-groups \(E\), and describes group cohomology as polynomial, allowing algebraic geometry in the algebraically closed-field setting of Benson--Carlson--Rickard. That polynomial description needs its prime convention: at odd primes the usual cohomology has exterior factors as well, or a polynomial description requires passage past nilpotents. The excerpt does not make this qualification, so its description is retained with that issue explicit.

For \(A(n)\) it says the quasi-elementary sub-Hopf algebras are exterior, but their Ext groups carry a bigrading. Conventionally the two gradings are cohomological and internal, as in \(\operatorname{Ext}^{s,t}(k,k)\); the exact objects and grading here are unspecified. The research task is to develop the appropriate bigraded support geometry and complete the intended classification, rather than discard one grading without justification. Hovey describes this as the obstacle then encountered and says something should be possible. The referenced Hovey--Palmieri abstract and DVI links have no supplied bibliographic identification.
\subsection{Short English statement}
Extend support-theoretic thick-subcategory classifications to the intended connected cocommutative Hopf algebras, especially \(A(n)\), using the reported quasi-elementary reduction and an appropriate bigraded Ext geometry. The category, tensor-ideal scope, prime and field assumptions need to be fixed.
\subsection{Sources}
\begin{itemize}
\item Local UnsolvedMath registry, permanent ID 11100032, source code catalogue entry; proposer Mark Hovey, ``Applications,'' item 8.
\item \emph{Hovey's Algebraic Topology Problem List}, surviving mirror: \url{https://www-users.cse.umn.edu/~tlawson/hovey/}. Recorded access 2026-07-29; extraction html-ordered-list. The mirror has 65 top-level items across nine category pages.
\item Benson--Carlson--Rickard are credited for the classification to extend; Hovey--Palmieri for partial success and reduction to quasi-elementary cases. Full papers, the abstract link and DVI link are not identified in the supplied record. No file was fetched.
\item Supplied status partially\_solved and PARTIAL-PROGRESS. The summary mentions broad stratification progress but says a clean bigraded-geometric classification for the stated case remains incomplete, without exact theorems. Current primary-literature status and redistribution terms remain NEEDS\_REVIEW. The registry assigns L4 whereas background records default L3.
\end{itemize}
\end{document}
